\chapter{Conclusion}
\label{ca:chap:conclusion}

In this report, we introduced provider-selective distillation, in which a small student
with a fixed budget chooses which priced, black-box model to learn from, and built a
testbed to study it. So far we find that text-only rewards are a safer basis than
cross-tokenizer logits, that price rather than provider strength decides value per unit of
budget, and that a selector which learns where progress happens matches the best single
provider without being told which one it is. Whether the best affordable provider differs
by domain or changes during training, and whether adaptivity can beat the best fixed
choice, remain open. Next semester we will measure marginal value with enough power to
settle H1 and H2, complete the H3 comparison at more budgets and against more selectors,
and explain the training instability.

\chapter*{Limitations}
\addcontentsline{toc}{chapter}{Limitations}
\label{ca:chap:limitations}

While the testbed lets us measure provider value directly, it simulates the setting rather
than reproducing it. Providers are open models that answer from a fixed cache of four
responses per prompt, prices are a compute-based proxy rather than real API prices, and
three of the five real providers share the student's model family. The student learns
almost only on Countdown at this scale, so most of our evidence comes from one domain. Three
seeds give wide intervals, and the branches saved only aggregate accuracies, so the planned
bootstrap over selection-set items still needs a re-evaluation of every branch. We also did
not vary the student's size, run more than one budget, or study multi-turn agentic tasks,
and the motivating experiment used a different teacher from the main provider pool.
